\chapter{Mérési környezet}

Jelen fejezet a kísérleti mérések elvégzéséhez kialakított hardver- és szoftverkörnyezetet mutatja be.
A fejezet célja, hogy az eredmények reprodukálhatóságához szükséges összes technikai részletet
rögzítse: az alkalmazott eszközök specifikációját, a szoftverkomponensek verzióit és a hálózati
topológiát.

%----------------------------------------------------------------------------
\section{Hardverkörnyezet}
%----------------------------------------------------------------------------

\subsection{Tesztplatform specifikációja}

% TODO: A mérésekhez használt eszköz(ök) pontos leírása (pl. Raspberry Pi 4B, ARM Cortex-M4
% alapú mikrovezérlő, x86 referenciagép stb.) — processzor, órajel, RAM, flash mérete.

\subsection{Hálózati topológia}

% TODO: A mérési elrendezés topológiájának leírása (kliens–broker–szerver kapcsolat,
% hálózati interfész típusa, lokális vs. konténeres kommunikáció).

\subsection{Energiaellátás és mérőeszközök}

% TODO: Ha energiamérés is történt: milyen eszközzel (pl. INA219 áramszenzor, PowerProfiler Kit),
% hogyan lett bekötve. Ha nem, ez az alfejezet elhagyható.

%----------------------------------------------------------------------------
\section{Szoftverkörnyezet}
%----------------------------------------------------------------------------

\subsection{Operációs rendszer és alapkönyvtárak}

% TODO: OS verzió, kernel, libc, gcc/clang verzió, cmake verzió.

\subsection{Kriptográfiai könyvtárak}

% TODO: wolfSSL / OpenSSL-OQS verzió, liboqs verzió, oqs-provider verzió, fordítási flagek
% (pl. -O2, konstans idejű fordítás, hardveres gyorsítás be/ki).

\subsection{MQTT infrastruktúra}

% TODO: Eclipse Mosquitto verzió, konfigurációs paraméterek (port, TLS verzió, cipher suite),
% kliens könyvtár (paho, mosquitto_pub/sub stb.).

\subsection{Konténerizáció és reprodukálhatóság}

% TODO: Docker image verziók, Dockerfile főbb lépései, volume- és portmapping.

%----------------------------------------------------------------------------
\section{A mérési környezet összefoglalása}
%----------------------------------------------------------------------------

% TODO: Rövid összefoglaló táblázat a főbb komponensekről és verziókról.
% Esetleg megjegyzés arról, hogy a környezet hogyan tér el a valódi IoT deploymenttől
% (pl. teljesítménybeli különbség egy gyári Cortex-M0 és a tesztplatform között).
